\chapter{المنتصفات والمتوازيات}\label{ch:g8:midpoints}

صِل بين منتصفَي ضلعين من مثلث: \omterm{def:g6:lines:objects}{فالقطعة} التي تحصل عليها
موازية دائمًا للضلع الثالث، وطولها نصف طوله بالضبط. وهذه
«مبرهنة المنتصفين» وعكسها أول مذاق لفكرة كبرى ---
أن المتوازيات تقطع الأطوال \omterm{def:g7:prop:table}{بتناسب} --- تتفتّح
في مبرهنة طاليس في \cref{ch:g9:thales}.

\section{مبرهنة المنتصفين}

\begin{theorem}[مبرهنة المنتصفين]\label{thm:g8:midpoints:direct}
في مثلث $ABC$، ليكن $I$ منتصف $[AB]$ و $J$
منتصف $[AC]$. عندئذٍ يكون \omterm{def:g6:lines:objects}{المستقيم} $(IJ)$ \omterm{def:g6:lines:perp}{موازيًا} \omterm{def:g6:lines:objects}{للمستقيم} $(BC)$، و
\[
IJ = \frac{BC}{2} .
\]
\end{theorem}

\begin{proof}[فكرة البرهان]
ليكن $K$ \omterm{def:g7:central:def}{مماثلة} $I$ بالنسبة إلى $J$ (أي نصف دورة،
\cref{ch:g7:central}). \omterm{def:g4:geometry:polygon}{فالرباعي} $AICK$ قطراه $[IK]$
و $[AC]$ يتقاطعان في منتصفهما المشترك $J$: إذن هو
\omterm{def:g7:central:parallelogram}{متوازي أضلاع}، ومنه $KC$ \omterm{def:g6:lines:perp}{موازٍ} \omterm{def:g6:lines:objects}{للمستقيم} $AI$، أي\ \omterm{def:g6:lines:objects}{للمستقيم} $(AB)$، و
$KC = AI = IB$. ثم إن $IBCK$ له ضلعان متقابلان ($[IB]$ و
$[KC]$) متوازيان ومتساويان: فهو \omterm{def:g7:central:parallelogram}{متوازي أضلاع} أيضًا
(\cref{thm:g7:central:recognize})، إذن $(IK) \parallel (BC)$ ---
أي $(IJ) \parallel (BC)$ --- و $BC = IK = 2\,IJ$.
\end{proof}

\begin{omfigure}
\begin{tikzpicture}[scale=1.1]
  \coordinate (A) at (1.7,2.9);
  \coordinate (B) at (0,0);
  \coordinate (C) at (4.6,0);
  \coordinate (I) at ($(A)!0.5!(B)$);
  \coordinate (J) at ($(A)!0.5!(C)$);
  \draw[thick] (A) -- (B) -- (C) -- cycle;
  \draw[omThm, very thick] (I) -- (J);
  \fill (A) circle (1.5pt) node[above, font=\small] {$A$};
  \fill (B) circle (1.5pt) node[below left, font=\small] {$B$};
  \fill (C) circle (1.5pt) node[below right, font=\small] {$C$};
  \fill[omThm] (I) circle (1.5pt) node[left, font=\small] {$I$};
  \fill[omThm] (J) circle (1.5pt) node[right, font=\small] {$J$};
  \draw ($(A)!0.22!(B)$) ++(-0.09,-0.05) -- ++(0.18,0.1);
  \draw ($(A)!0.72!(B)$) ++(-0.09,-0.05) -- ++(0.18,0.1);
  \draw ($(A)!0.25!(C)$) ++(-0.06,-0.1) -- ++(0.12,0.2);
  \draw ($(A)!0.25!(C)$) ++(0.02,-0.1) -- ++(0.12,0.2);
  \draw ($(A)!0.75!(C)$) ++(-0.06,-0.1) -- ++(0.12,0.2);
  \draw ($(A)!0.75!(C)$) ++(0.02,-0.1) -- ++(0.12,0.2);
\end{tikzpicture}

{\small \omterm{def:g6:lines:objects}{قطعة} المنتصفين $[IJ]$ (بالأحمر) موازية للضلع
الثالث $(BC)$ وطولها نصف طوله.}
\end{omfigure}

\begin{example}\label{ex:g8:midpoints:direct}
في مثلث $ABC$ فيه $BC = 9$ سم، يُوصَل بين منتصف $I$ \omterm{def:g6:lines:objects}{للقطعة} $[AB]$ و
منتصف $J$ \omterm{def:g6:lines:objects}{للقطعة} $[AC]$. ودون أيّ قياس:
$(IJ) \parallel (BC)$ و $IJ = \frac92 = 4.5$ سم. والمثلث الصغير
$AIJ$ نسخة بنصف \omterm{def:g6:measure:volume}{الحجم} من $ABC$ --- ومحيطه نصف
\omterm{def:g6:measure:perimeter}{محيط} $ABC$ أيضًا.
\end{example}

\begin{theorem}[العكس]\label{thm:g8:midpoints:converse}
في مثلث $ABC$، \omterm{def:g6:lines:objects}{المستقيم} المارّ بمنتصف $I$ \omterm{def:g6:lines:objects}{للقطعة} $[AB]$ و
\emph{\omterm{def:g6:lines:perp}{الموازي}} \omterm{def:g6:lines:objects}{للمستقيم} $(BC)$ يقطع الضلع $[AC]$ في منتصفه.
\end{theorem}
\admitted

\begin{example}\label{ex:g8:midpoints:converse}
مسار يعبر حقلًا مثلثيًّا، منطلقًا من وسط حافة
وسائرًا \omterm{def:g6:lines:perp}{موازيًا} للحافة المقابلة. فالعكس يضمن
أن المسار يخرج بالضبط من وسط الحافة الأخرى --- بلا
قياس في الجهة البعيدة.
\end{example}

\begin{method}[الاختيار بين المبرهنة وعكسها]\label{met:g8:midpoints:choose}
\begin{enumerate}
  \item \emph{منتصفان معلومان} $\Rightarrow$ اِستعمل المبرهنة
        المباشرة: واِستنتج التوازي ونصف الطول.
  \item \emph{منتصف واحد ومتوازٍ معلومان} $\Rightarrow$ اِستعمل
        العكس: واِستنتج أن نقطة التقاطع منتصف.
  \item واُكتب أيّ مثلث وأيّ منتصفات وأيّ مبرهنة
        تستعمل --- جملة لكل منها.
\end{enumerate}
\end{method}

\section{مثلثات بنصف الحجم}

\begin{proposition}[مثلث المنتصفات]\label{prop:g8:midpoints:medial}
الوصل بين منتصفات أضلاع مثلث الثلاثة يقطعه إلى
أربعة مثلثات متساوية المساحات، كل منها نسخة بنصف \omterm{def:g6:measure:volume}{الحجم} من الأصل.
\end{proposition}

\begin{proof}
كل ضلع من مثلث المنتصفات \omterm{def:g6:lines:perp}{موازٍ} لضلع من $ABC$ و
نصفه طولًا (\cref{thm:g8:midpoints:direct}، مطبَّقة ثلاث مرات).
وللمثلثات الصغيرة الأربعة أضلاع بالأطوال الثلاثة نفسها، إذن
هي نسخ متطابقة؛ وهي معًا تبلّط $ABC$، إذن لكل منها
رُبع مساحته.
\end{proof}

\begin{omfigure}
\begin{tikzpicture}[scale=1.05]
  \coordinate (A) at (1.9,3.0);
  \coordinate (B) at (0,0);
  \coordinate (C) at (5.0,0);
  \coordinate (I) at ($(A)!0.5!(B)$);
  \coordinate (J) at ($(A)!0.5!(C)$);
  \coordinate (K) at ($(B)!0.5!(C)$);
  \draw[thick] (A) -- (B) -- (C) -- cycle;
  \fill[omThm!12] (I) -- (J) -- (K) -- cycle;
  \draw[omThm, thick] (I) -- (J) -- (K) -- cycle;
  \fill (A) circle (1.4pt) node[above, font=\small] {$A$};
  \fill (B) circle (1.4pt) node[below left, font=\small] {$B$};
  \fill (C) circle (1.4pt) node[below right, font=\small] {$C$};
  \fill[omThm] (I) circle (1.4pt) node[left, font=\small] {$I$};
  \fill[omThm] (J) circle (1.4pt) node[right, font=\small] {$J$};
  \fill[omThm] (K) circle (1.4pt) node[below, font=\small] {$K$};
\end{tikzpicture}

{\small مثلث المنتصفات $IJK$ يقطع $ABC$ إلى أربعة مثلثات
متساوية --- وهي صورة تستحقّ التذكّر.}
\end{omfigure}

\begin{remark}[نحو طاليس]
مبرهنة المنتصفين هي حالة «النصف» من واقعة أعمّ:
\omterm{def:g6:lines:objects}{فالمستقيم} \omterm{def:g6:lines:perp}{الموازي} لضلع من مثلث يقطع الضلعين الآخرين
\emph{بنسب متساوية} --- ثلثًا وثلثًا، وخُمسين
وخُمسين\,\dots\ وتلك العبارة العامّة هي مبرهنة طاليس
(\cref{ch:g9:thales})؛ وحالة المنتصفين هي الوحيدة المحتاج إليها هذه
السنة.
\end{remark}

\section{تمارين}

\begin{exercise}[$\star$]\label{exo:g8:midpoints:1}
في مثلث $RST$، $M$ منتصف $[RS]$ و $N$
منتصف $[RT]$، مع $ST = 7$ سم. فماذا يمكن أن يقال عن \omterm{def:g6:lines:objects}{المستقيم}
$(MN)$ والطول $MN$؟ اُذكر المبرهنة المستعملة.
\end{exercise}

\begin{exercise}[$\star$]\label{exo:g8:midpoints:2}
في مثلث $ABC$، $I$ منتصف $[AB]$ \omterm{def:g6:lines:objects}{والمستقيم}
المارّ بالنقطة $I$ \omterm{def:g6:lines:perp}{الموازي} \omterm{def:g6:lines:objects}{للمستقيم} $(BC)$ يقطع $[AC]$ في $J$، مع
$AC = 11$ سم. اِحسب $AJ$، مع ذكر المبرهنة المستعملة.
\end{exercise}

\begin{exercise}[$\star$]\label{exo:g8:midpoints:3}
اُرسم مثلثًا $ABC$ حيث $AB = 8$ سم، و $AC = 6$ سم، و $BC = 7$ سم،
وضع منتصف $I$ \omterm{def:g6:lines:objects}{للقطعة} $[AB]$ ومنتصف $J$ \omterm{def:g6:lines:objects}{للقطعة} $[AC]$، وقِس $IJ$
للتحقّق من المبرهنة. واِحسب \omterm{def:g6:measure:perimeter}{محيط} $AIJ$ دون
قياس أيّ شيء آخر.
\end{exercise}

\begin{exercise}[$\star$]\label{exo:g8:midpoints:4}
المثلث $IJK$ هو مثلث المنتصفات للمثلث $ABC$، الذي تقيس أضلاعه $10$
و $12$ و $16$ سم. أعطِ أضلاع $IJK$ الثلاثة ومحيطه.
فكيف يتقارن المحيطان؟
\end{exercise}

\begin{exercise}[$\star$]\label{exo:g8:midpoints:5}
\omterm{def:g6:measure:area}{مساحة} مثلث $36$ سم$^2$. فما \omterm{def:g6:measure:area}{مساحة} مثلث
منتصفاته؟ وما \omterm{def:g6:measure:area}{مساحة} كل واحد من المثلثات الصغيرة الأربعة؟
\end{exercise}

\begin{exercise}[$\star\star$]\label{exo:g8:midpoints:6}
في مثلث $ABC$، $I$ منتصف $[AB]$، و $J$ منتصف
$[AC]$، و $K$ منتصف $[BC]$. بيّن أن $IJKB$ \omterm{def:g7:central:parallelogram}{متوازي أضلاع}.
(قارن $[IJ]$ و $[BK]$: أمتوازيان؟ أمتساويان؟)
\end{exercise}

\begin{exercise}[$\star\star$]\label{exo:g8:midpoints:7}
شراع مثلثي حافته السفلى طولها $6.4$ م. ودرزة تقوية
تصل بين منتصفَي الحافتين الأخريين. فما طول الدرزة
اللازم؟ وماذا لو وصلت الدرزة بدل ذلك بين النقطتين الواقعتين على
رُبع كل حافة، اِنطلاقًا من الرأس الأعلى؟ (خمّن
برسم؛ والبرهان لطاليس، \cref{ch:g9:thales}.)
\end{exercise}

\begin{exercise}[$\star\star$]\label{exo:g8:midpoints:8}
في المثلث $ABC$، \omterm{def:g6:shapes:triangles}{القائم الزاوية} في $A$، $M$ منتصف
الوتر $[BC]$، \omterm{def:g6:lines:perp}{والموازي} \omterm{def:g6:lines:objects}{للمستقيم} $(AB)$ المارّ بالنقطة $M$ يلتقي
$[AC]$ في $N$.
\begin{enumerate}
  \item بيّن أن $N$ منتصف $[AC]$.
  \item واِشرح لماذا يكون $(MN)$ عموديًّا على $(AC)$.
\end{enumerate}
\end{exercise}

\begin{exercise}[$\star\star\star$]\label{exo:g8:midpoints:9}
ليكن $ABCD$ رباعيًا \emph{كيفما كان}، و $I$ و $J$ و $K$ و $L$
منتصفات $[AB]$ و $[BC]$ و $[CD]$ و $[DA]$ (وهو تخمين
\cref{exo:g7:central:11}).
\begin{enumerate}
  \item طبّق مبرهنة المنتصفين في المثلث $ABC$ على
        \omterm{def:g6:lines:objects}{القطعة} $[IJ]$، وفي المثلث $ACD$ على \omterm{def:g6:lines:objects}{القطعة}
        $[LK]$: وقارن كلًّا منهما \omterm{def:g4:geometry:circle}{بالقطر} $[AC]$.
  \item واِستنتج أن $IJKL$ \omterm{def:g7:central:parallelogram}{متوازي أضلاع} دائمًا.
\end{enumerate}
\end{exercise}

\section{مسألة: المتوسطات الثلاثة ومركز الثقل}

\begin{problem}[{مسألة نهاية الأسبوع --- متوسطات المثلث تلتقي
في نقطة واحدة، على ثلثَي كل واحد منها}]
\label{pb:g8:midpoints:1}
اِقطع مثلثًا من ورق مقوّى فسيتوازن، مسطّحًا تمامًا،
على سنّ قلم موضوع في نقطة خاصّة واحدة: هي
\emph{مركز ثقل}\index{مركز ثقل} المثلث.
وتجد هذه المسألة تلك النقطة. و\emph{المتوسط}\index{متوسط} في
مثلث هو \omterm{def:g6:lines:objects}{القطعة} الواصلة بين رأس و\emph{منتصف}
الضلع المقابل. وستبرهن --- بمبرهنة المنتصفين
مستعملةً مرتين، بطريقة غير متوقّعة --- على أن المتوسطات الثلاثة
تمرّ كلها بنقطة واحدة، وستحدّد موضعها بدقّة.

\medskip
\noindent\textbf{الجزء الأول --- إحماء.}
\begin{enumerate}
  \item اُرسم مثلثًا كبيرًا $ABC$ (بلا شكل خاصّ)، واِبنِ
        منتصف $I$ \omterm{def:g6:lines:objects}{للقطعة} $[AB]$، ومنتصف $J$ \omterm{def:g6:lines:objects}{للقطعة} $[AC]$، ومنتصف $K$
        \omterm{def:g6:lines:objects}{للقطعة} $[BC]$، واُرسم المتوسطات الثلاثة $[AK]$ و $[BJ]$ و $[CI]$.
        فماذا تلاحظ؟
  \item بَرهِن على أن المتوسط يقطع المثلث إلى مثلثين
        متساويي \omterm{def:g6:measure:area}{المساحة}. (قارن القواعد والارتفاعات، واِستعمل
        صيغة \omterm{def:g6:measure:area}{المساحة}: القاعدة $\times$ الارتفاع $\div\ 2$.) ومن أجل
        مثلث مساحته $36$ سم$^2$، ما مساحتا
        \omterm{def:g6:lines:objects}{القطعتين}؟
  \item وفي المثلث $ABC$، ماذا تقول
        \cref{thm:g8:midpoints:direct} عن \omterm{def:g6:lines:objects}{القطعة}
        $[IJ]$؟
\end{enumerate}

\medskip
\noindent\textbf{الجزء الثاني --- متوسطان يلتقيان على ثلثَي
الطريق.}
يتقاطع المتوسطان $[BJ]$ و $[CI]$ في نقطة؛ سمِّها $G$. وليكن
$M$ منتصف $[GB]$ و $N$ منتصف $[GC]$.
\begin{enumerate}[resume]
  \item طبّق مبرهنة المنتصفين \emph{في المثلث $GBC$}:
        فماذا تقول عن \omterm{def:g6:lines:objects}{القطعة} $[MN]$؟
  \item اِستنتج أن $(IJ) \parallel (MN)$ و $IJ = MN$، و
        اُختم حسب \cref{thm:g7:central:recognize} بأن $IJNM$
        \omterm{def:g7:central:parallelogram}{متوازي أضلاع}.
  \item اِشرح لماذا تقع $I$ و $N$ كلتاهما على المتوسط $(CI)$،
        و $J$ و $M$ كلتاهما على المتوسط $(BJ)$ --- فيكون
        قطرا \omterm{def:g7:central:parallelogram}{متوازي الأضلاع} $IJNM$ هما $[IN]$ و
        $[JM]$، ويتقاطعان في $G$ بالضبط. فماذا يقول
        \emph{تعريف} \omterm{def:g7:central:parallelogram}{متوازي الأضلاع} حينئذٍ عن $G$؟
  \item اِستنتج أن $BM = MG = GJ$، واُختم:
        \[
        BG = 2 \times GJ
        \]
        --- فالنقطة $G$ تقع على المتوسط $[BJ]$ على ثلثَي
        الطريق من الرأس $B$. واُذكر النتيجة المشابهة على
        المتوسط $[CI]$ وبَرِّرها.
  \item والآن اِنسَ $[CI]$ وشغّل الاستدلال نفسه بالضبط على
        زوج المتوسطين $[BJ]$ و $[AK]$: فنقطة تقاطعهما
        تقع أيضًا على ثلثَي $[BJ]$ من $B$. اِشرح لماذا
        يفرض هذا أن تكون النقطة $G$ \emph{نفسها} --- و
        اِستنتج أن المتوسطات الثلاثة كلها تمرّ بالنقطة $G$.
  \item في مثلث يقيس فيه المتوسط $[BJ]$ طولًا قدره $9$ سم،
        كم تبعد $G$ عن $B$، وعن $J$؟
  \item بَرِّر (في جملة) أن $AG = 2 \times GK$ أيضًا.
\end{enumerate}

\medskip
\noindent\textbf{الجزء الثالث --- ستّ شرائح متساوية، وإحداثيات.}
المتوسطات الثلاثة تقطع المثلث إلى ستة مثلثات صغيرة
حول $G$.
\begin{enumerate}[resume]
  \item في المثلث $GBC$، \omterm{def:g6:lines:objects}{القطعة} $[GK]$ متوسط.
        فاستنتج أن المثلثين الصغيرين $GBK$ و $GKC$ لهما
        مساحتان متساويتان.
  \item وباستعمال السؤال 2 في المثلثين $ABK$ و $AKC$ (وكلاهما
        مقطوع بالمتوسط $[AK]$ للمثلث $ABC$)، بيّن أن
        المثلثين $ABG$ و $ACG$ متساويا \omterm{def:g6:measure:area}{المساحة}. وبالتكرار
        مع متوسط آخر، اُختم بأن المثلثات الثلاثة
        $ABG$ و $BCG$ و $CAG$ \omterm{def:g6:measure:area}{مساحة} كل منها ثلث $ABC$.
  \item وكل واحد من هذه المثلثات الثلاثة تقطعه \omterm{def:g6:lines:objects}{قطعة}
        متوسط إلى نصفين (\omterm{def:g6:lines:objects}{فالقطعة} $[GI]$ مثلًا متوسط في
        المثلث $ABG$ --- منظورًا إليه من أيّ رأس؟). فاستنتج: الشرائح الست
        الصغيرة حول $G$ لها \omterm{def:g6:measure:area}{المساحة} نفسها كلها، وهي سُدس
        المثلث.
  \item على شبكة إحداثيات، ضع $A(0, 0)$ و $B(6, 0)$ و
        $C(0, 6)$، مع $K(3, 3)$ منتصف $[BC]$. واِستعمل
        مبرهنة الثلثين في الجزء الثاني لتحسب إحداثيَي
        $G$ على المتوسط $[AK]$؛ ثم تحقّق من أن النقطة
        نفسها تقع على ثلثَي المتوسط $[BJ]$، حيث
        $J(0, 3)$. وقارن إحداثيَي $G$ بالمتوسطين الحسابيين
        \[
        \frac{0 + 6 + 0}{3},
        \qquad
        \frac{0 + 0 + 6}{3} .
        \]
  \item خاتمة التوازن: باستعمال السؤالين 2 و8، اِشرح لماذا
        يتوازن المثلث المقوّى على حدّ سكين موضوع
        على أيّ متوسط --- بمقادير متساوية من الورق المقوّى في كل
        جهة --- ولماذا يجب إذن أن توضع سنّ القلم
        عند $G$، وهي النقطة التي يسمّيها الفيزيائيون \omterm{pb:g8:midpoints:1}{مركز
        الثقل}. (والبرهان الصارم تمامًا على التوازن يحتاج
        حساب التكامل في مجلّدات الجامعة؛ وتساوي
        المساحات يجعله معقولًا اليوم.)
\end{enumerate}
\end{problem}
